% !TEX program = pdflatex
\documentclass[11pt,a4paper]{article}
\usepackage[margin=25mm]{geometry}
\usepackage{amsmath}
\title{A Small Study of Writing Workflows}
\author{Catalex Demo Project}
\date{}
\begin{document}
\maketitle
\begin{abstract}
This sample paper provides three pages of text, equations, and a table for
exploring the PDF workspace. All results are invented for demonstration.
\end{abstract}
\section{Introduction}
Writing involves moving between ideas, source text, and the finished document.
A useful workspace keeps those views close without competing for attention.
This example explores how an author might compare several revision workflows.

The central question is whether a compact conversation helps the author stay
focused, and when a larger conversation becomes useful. Neither layout is
expected to fit every task. The author should be able to switch layouts while
keeping the same passage in view.
\subsection{Study questions}
\begin{enumerate}
\item How often do authors switch between source and preview?
\item Does a larger conversation help with longer explanations?
\item Can readers return to the same passage after changing the layout?
\end{enumerate}
\paragraph{Reading exercise.}
Scroll to the second page, then switch between Split and PDF or between Half
and Full chat. Use Fit width or Fit height to compare the two fitting modes.

\newpage
\section{Method}
Each participant revises two short passages using both workspace layouts.
The order is counterbalanced, and each passage has the same target length.
Participants can inspect the PDF at any time while editing the source.

We summarize the revision cost as
\begin{equation}
C = \alpha T + \beta S + \gamma E,
\end{equation}
where $T$ is task time, $S$ is the number of view switches, and $E$ counts
remaining errors. The weights express the relative importance of each term;
they are not measurements from a real study.

\subsection{Procedure}
\begin{enumerate}
\item Read the passage and identify its main argument.
\item Ask for a clearer introduction without changing the technical claims.
\item Review the suggestion and keep only useful changes.
\item Inspect the PDF and rate the ease of recovering the reading position.
\end{enumerate}
\subsection{Observations}
The source and preview support different kinds of inspection. Source text
helps with precise edits, while the PDF exposes page breaks, spacing, and the
overall shape of the document. A stable preview makes those transitions easier
to follow.

\newpage
\section{Illustrative Results}
\begin{center}
\begin{tabular}{lrrr}
\hline
Workflow & Time (min) & View switches & Rating (1--7) \\
\hline
Compact chat & 12.4 & 18 & 5.8 \\
Expanded chat & 11.7 & 14 & 6.1 \\
Flexible layout & 10.9 & 12 & 6.4 \\
\hline
\end{tabular}
\end{center}
These invented values simply provide a table to examine at different zoom
levels. They do not establish that one workflow is better than another.

\section{Discussion}
Compact and expanded layouts can serve complementary purposes. A short
question may fit naturally beside the draft, while a long explanation may
benefit from more room. The transition should help the author follow the
conversation without pulling attention toward the page resize.

\section{Conclusion}
A predictable workspace preserves context through ordinary layout changes.
The same source remains editable, the PDF remains readable, and the author
decides which view is most useful for the next step.
\end{document}
